% GENERATED FILE. Edit canonical lessons, then run npm run generate:books.
% canonical-source-hash: fnv1a64:38cf408112df5180
% canonical-lessons: AR-C441-istiqbal, AR-C441-nazil, AR-C441-iqama, AR-C441-ghadara, AR-C441-akhla

\chapter{At reception}
\label{ch:ar-at-reception}
\hlchaptermodality{441}

\begin{chapteropening}
\textbf{By the end of this chapter:} \emph{I can say reception, a hotel guest, a residence, to depart from and to vacate.}

Complete the last lesson of chapter 441: I can say reception, a hotel guest, a residence, to depart from and to vacate.
\end{chapteropening}

\section[\texorpdfstring{\ar{استقبال} (istiqbāl)}{استقبال (istiqbāl)}]{\ar{استقبال} (istiqbāl) — reception; a welcome}

\label{lesson:AR-C441-istiqbal}



\begin{quote}
Before the new one: say the Arabic for to accompany, then the Arabic for to give a lift.
\end{quote}

\subsection*{You'll want to know: \ar{استقبال}}
\textbf{\ar{استقبال}} — \emph{istiqbāl} — \textquotedblleft{}reception; a welcome\textquotedblright{}.

\begin{quote}
\ar{أعطاني} \ar{موظف} \ar{الاستقبال} \ar{المفتاح}.
\end{quote}

\begin{quote}
\emph{aʿṭānī muwaẓẓaf al-istiqbāl al-miftāḥ.}
\end{quote}

\begin{quote}
The receptionist gave me the key.
\end{quote}

\begin{cousinweb}
\textbf{\ar{استقبال}} (\emph{istiqbāl}) is built on the native Arabic root \emph{q-b-l}, \textquotedblleft{}to be in front of; to accept\textquotedblright{} (Lane). It is the verbal noun of Form X, \textquotedblleft{}going out to meet\textquotedblright{}.
\end{cousinweb}

\begin{grammarlens}[title={What you've built}]
One more word for getting around.
\end{grammarlens}

\subsection*{Your turn}
\begin{itemize}
\raggedright
  \item \emph{Say it:} \emph{istiqbāl}
  \item \emph{Say it:} \emph{istiqbāl}, once more
  \item \emph{Say it:} \emph{aʿṭānī muwaẓẓaf al-istiqbāl al-miftāḥ.} — The receptionist gave me the key.
  \item \emph{Recall:} say the Arabic for to accompany, then the Arabic for to give a lift, then say \emph{istiqbāl} again
\end{itemize}

\begin{tcolorbox}[breakable,colback=teal!4,colframe=teal!35!black,title={Before you move on}]
What does \ar{استقبال} mean? (\textquotedblleft{}Reception; a welcome\textquotedblright{}.) Say it once more.
\end{tcolorbox}
\section[\texorpdfstring{\ar{نزيل} (nazīl)}{نزيل (nazīl)}]{\ar{نزيل} (nazīl) — a hotel guest; a lodger}

\label{lesson:AR-C441-nazil}



\begin{quote}
Before the new one: say the Arabic for to give a lift, then the Arabic for reception.
\end{quote}

\subsection*{You'll want to know: \ar{نزيل}}
\textbf{\ar{نزيل}} — \emph{nazīl} — \textquotedblleft{}a hotel guest; a lodger\textquotedblright{}.

\begin{quote}
\ar{طلب} \ar{النزيل} \ar{غرفة} \ar{بعيدة} \ar{عن} \ar{الشارع}.
\end{quote}

\begin{quote}
\emph{ṭalaba an-nazīl ghurfa baʿīda ʿan ash-shāriʿ.}
\end{quote}

\begin{quote}
The guest asked for a room far from the street.
\end{quote}

Plural: \ar{نزلاء}, \emph{nuzalāʾ}.

\begin{cousinweb}
\textbf{\ar{نزيل}} (\emph{nazīl}) is built on the native Arabic root \emph{n-z-l}, \textquotedblleft{}to alight, to stop at a place\textquotedblright{} (Lane). A \emph{nazīl} is one who has stopped somewhere for the night.
\end{cousinweb}

\begin{grammarlens}[title={What you've built}]
One more word for getting around.
\end{grammarlens}

\subsection*{Your turn}
\begin{itemize}
\raggedright
  \item \emph{Say it:} \emph{nazīl}
  \item \emph{Say it:} \emph{nazīl}, once more
  \item \emph{Say it:} \emph{ṭalaba an-nazīl ghurfa baʿīda ʿan ash-shāriʿ.} — The guest asked for a room far from the street.
  \item \emph{Recall:} say the Arabic for to give a lift, then the Arabic for reception, then say \emph{nazīl} again
\end{itemize}

\begin{tcolorbox}[breakable,colback=teal!4,colframe=teal!35!black,title={Before you move on}]
What does \ar{نزيل} mean? (\textquotedblleft{}A hotel guest; a lodger\textquotedblright{}.) Say it once more.
\end{tcolorbox}
\section[\texorpdfstring{\ar{إقامة} (iqāma)}{إقامة (iqāma)}]{\ar{إقامة} (iqāma) — a residence; a stay}

\label{lesson:AR-C441-iqama}



\begin{quote}
Before the new one: say the Arabic for reception, then the Arabic for a hotel guest.
\end{quote}

\subsection*{You'll want to know: \ar{إقامة}}
\textbf{\ar{إقامة}} — \emph{iqāma} — \textquotedblleft{}a residence; a stay\textquotedblright{}.

\begin{quote}
\ar{كانت} \ar{إقامتنا} \ar{في} \ar{الفندق} \ar{مريحة}.
\end{quote}

\begin{quote}
\emph{kānat iqāmatunā fī al-funduq murīḥa.}
\end{quote}

\begin{quote}
Our stay in the hotel was comfortable.
\end{quote}

\emph{biṭāqat iqāma} is a residence permit.

\begin{cousinweb}
\textbf{\ar{إقامة}} (\emph{iqāma}) is built on the native Arabic root \emph{q-w-m}, \textquotedblleft{}to stand\textquotedblright{} (Lane). It is the verbal noun of Form IV, \textquotedblleft{}staying on in a place\textquotedblright{}.
\end{cousinweb}

\begin{grammarlens}[title={What you've built}]
One more word for getting around.
\end{grammarlens}

\subsection*{Your turn}
\begin{itemize}
\raggedright
  \item \emph{Say it:} \emph{iqāma}
  \item \emph{Say it:} \emph{iqāma}, once more
  \item \emph{Say it:} \emph{kānat iqāmatunā fī al-funduq murīḥa.} — Our stay in the hotel was comfortable.
  \item \emph{Recall:} say the Arabic for reception, then the Arabic for a hotel guest, then say \emph{iqāma} again
\end{itemize}

\begin{tcolorbox}[breakable,colback=teal!4,colframe=teal!35!black,title={Before you move on}]
What does \ar{إقامة} mean? (\textquotedblleft{}A residence; a stay\textquotedblright{}.) Say it once more.
\end{tcolorbox}
\section[\texorpdfstring{\ar{غادر} (ghādara)}{غادر (ghādara)}]{\ar{غادر} (ghādara) — to depart from; to leave (a place)}

\label{lesson:AR-C441-ghadara}



\begin{quote}
Before the new one: say the Arabic for a hotel guest, then the Arabic for a residence.
\end{quote}

\subsection*{You'll want to know: \ar{غادر}}
\textbf{\ar{غادر}} — \emph{ghādara} — \textquotedblleft{}to depart from; to leave (a place)\textquotedblright{}.

\begin{quote}
\ar{غادرنا} \ar{الفندق} \ar{في} \ar{الصباح}.
\end{quote}

\begin{quote}
\emph{ghādarnā al-funduq fī aṣ-ṣabāḥ.}
\end{quote}

\begin{quote}
We left the hotel in the morning.
\end{quote}

Present: \ar{يغادر}, \emph{yughādir}.

\begin{cousinweb}
\textbf{\ar{غادر}} (\emph{ghādara}) is built on the native Arabic root \emph{gh-d-r}, \textquotedblleft{}to leave behind\textquotedblright{} (Lane). It is Form III; the plain verb \emph{ghadara} means to betray.
\end{cousinweb}

\begin{grammarlens}[title={What you've built}]
One more word for getting around.
\end{grammarlens}

\subsection*{Your turn}
\begin{itemize}
\raggedright
  \item \emph{Say it:} \emph{ghādara}
  \item \emph{Say it:} \emph{ghādara}, once more
  \item \emph{Say it:} \emph{ghādarnā al-funduq fī aṣ-ṣabāḥ.} — We left the hotel in the morning.
  \item \emph{Recall:} say the Arabic for a hotel guest, then the Arabic for a residence, then say \emph{ghādara} again
\end{itemize}

\begin{tcolorbox}[breakable,colback=teal!4,colframe=teal!35!black,title={Before you move on}]
What does \ar{غادر} mean? (\textquotedblleft{}To depart from; to leave (a place)\textquotedblright{}.) Say it once more.
\end{tcolorbox}
\section[\texorpdfstring{\ar{أخلى} (akhlā)}{أخلى (akhlā)}]{\ar{أخلى} (akhlā) — to vacate; to evacuate}

\label{lesson:AR-C441-akhla}



\begin{quote}
Before the new one: say the Arabic for a residence, then the Arabic for to depart from.
\end{quote}

\subsection*{You'll want to know: \ar{أخلى}}
\textbf{\ar{أخلى}} — \emph{akhlā} — \textquotedblleft{}to vacate; to evacuate\textquotedblright{}.

\begin{quote}
\ar{أخلينا} \ar{الغرفة} \ar{قبل} \ar{الظهر}.
\end{quote}

\begin{quote}
\emph{akhlaynā al-ghurfa qabla aẓ-ẓuhr.}
\end{quote}

\begin{quote}
We vacated the room before noon.
\end{quote}

Present: \ar{يخلي}, \emph{yukhlī}.

\begin{cousinweb}
\textbf{\ar{أخلى}} (\emph{akhlā}) is built on the native Arabic root \emph{kh-l-w}, \textquotedblleft{}to be empty, to be alone\textquotedblright{} (Lane). It is Form IV, \textquotedblleft{}to make empty\textquotedblright{}.
\end{cousinweb}

\begin{grammarlens}[title={What you've built}]
One more word for getting around.
\end{grammarlens}

\subsection*{Your turn}
\begin{itemize}
\raggedright
  \item \emph{Say it:} \emph{akhlā}
  \item \emph{Say it:} \emph{akhlā}, once more
  \item \emph{Say it:} \emph{akhlaynā al-ghurfa qabla aẓ-ẓuhr.} — We vacated the room before noon.
  \item \emph{Recall:} say the Arabic for a residence, then the Arabic for to depart from, then say \emph{akhlā} again
\end{itemize}

\begin{tcolorbox}[breakable,colback=teal!4,colframe=teal!35!black,title={Before you move on}]
What does \ar{أخلى} mean? (\textquotedblleft{}To vacate; to evacuate\textquotedblright{}.) Say it once more.
\end{tcolorbox}
